% TEMPLATE-MILC-v3.tex - plantilla maestra UNICA de las guias MILC v3.
%
% La usan las cuatro familias de guias, cada una con su driver:
%   · Grados 6-11          -> scripts/build-guias-g11.py
%   · Bebras               -> scripts/build-guias-bebras.py
%   · Semillero            -> scripts/build-guias-semillero.py
%   · Territorio interior  -> scripts/build-guias-territorio-interior.py
%
% Antes habia tres copias casi identicas (~400 lineas iguales, 6 distintas):
% cada mejora habia que aplicarla tres veces, y en la practica no se hacia
% (el motor de recursos vivia solo en dos de ellas). Lo que varia entre
% familias esta parametrizado en 6 placeholders que cada driver llena
% (se nombran aqui SIN su sintaxis real: el builder reemplaza por texto plano,
% tambien dentro de los comentarios, y un valor multilinea rompe el preambulo):
%
%   HEADER_IZQ        encabezado izquierdo de pagina
%   HEADER_DER        encabezado derecho de pagina
%   PORTADA_ETIQUETA  rotulo sobre el numero grande (GRADO/MOMENTO/MODULO)
%   PORTADA_SERIE     linea de serie de la portada
%   PORTADA_ID        identificador de la guia en la portada
%   PORTADA_CIRCULO   el circulito de grado (°), o vacio si no aplica
%   PDF_TITULO        titulo en texto plano (sin \textbf ni comillas TeX) para
%                     los metadatos del PDF (/Title); lo produce lib_xelatex.texto_pdf
%
% Composicion (auditoria 2026-09-03): las bandas de estacion piden espacio
% (\Needspace*) y prohiben el salto justo despues (\nopagebreak), asi no quedan
% huerfanas al pie; las cajas largas (softbox, drawbox, infoband, puentebox) son
% `breakable`, asi una caja de media pagina ya no arrastra todo a la siguiente
% dejando la anterior al 40 %. Todo texto sobre mostaza o verde va en milcNegro
% (blanco sobre #E5B400 da 1,93:1 y sobre #58A923 2,95:1; ninguno pasa WCAG AA).
% El resto de claves las documenta content/guias/_SCHEMA.md. El script valida
% que no queden placeholders sin reemplazar antes de escribir el archivo final.

% Etiquetado PDF (latex-lab): /Lang, estructura y texto alternativo de las
% figuras, para lectores de pantalla. Debe ir ANTES de \documentclass.
\DocumentMetadata{lang=es-CO,tagging=on}
\documentclass[11pt,letterpaper]{article}
% headheight=24pt: la cabecera derecha lleva dos lineas (identidad de la guia
% y estacion actual). headsep se acorta en la misma medida para que el bloque
% de texto quede exactamente donde estaba con headheight=12pt.
\usepackage[margin=2.0cm,headheight=24pt,headsep=13pt]{geometry}
\usepackage[table]{xcolor}
\usepackage{fontspec}
\usepackage[spanish]{babel}
\usepackage{tikz}
% Librerías TikZ para los diagramas de las guías (bloque `recursos:`):
%  · babel  → babel en español vuelve activos caracteres como «>», y TikZ los usa
%             en sus opciones (>=stealth, ->). Sin esta librería, cualquier
%             diagrama con flechas revienta con «\language@active@arg> has an
%             extra }». Debe ir ANTES de que se dibuje nada.
%  · positioning        → sintaxis «below left=2pt and 3pt».
%  · decorations.pathreplacing → llaves ({brace}) para señalar rangos.
\usetikzlibrary{babel,positioning,decorations.pathreplacing}
\usepackage{graphicx}
% Circuitos: el trabajo de electrónica se hace hoy con micro:bit y sus sensores
% integrados (MakeCode, sin cableado), así que no se necesitan esquemáticos.
% Si algún día entran montajes con protoboard, basta descomentar esta línea:
% los diagramas .tex/.tikz declarados en `recursos:` ya se incrustan solos.
% \usepackage{circuitikz}
\usepackage[most]{tcolorbox}
\usepackage{tabularx}
\usepackage{array}
\usepackage{fancyhdr}
\usepackage{titlesec}
\usepackage[document]{ragged2e}  % bandera a la derecha en todo el cuerpo
\usepackage{enumitem}
\usepackage{microtype}
\usepackage{etoolbox}
\usepackage{needspace}
% hyperref al final del preambulo: metadatos del PDF (titulo, autor, idioma,
% familia), marcadores por estacion (\pdfbookmark) y enlaces sin recuadro.
\usepackage[hidelinks,
  pdftitle={El trabajo y la dignidad: el mismo oficio, tres maneras de vivirlo},
  pdfauthor={PhD. Álvaro Cárdenas Orozco},
  pdfsubject={Territorio interior · Educación socioemocional · Grado 11° · Momento 4 · MILC v3},
  pdflang={es-CO},
  bookmarksopen=true,bookmarksnumbered=false]{hyperref}

% Helvetica Neue no trae la flecha U+2192, asi que cada «→» escrito literalmente
% en el YAML se compilaba a NADA, en silencio (462 flechas en 61 guias: rutas del
% tipo «paleta → tipografia → lockup» salian rotas). Mapeamos el caracter a su
% version matematica, que si tiene glifo. Asi el autor puede seguir escribiendo
% «→» comodamente en el YAML.
\usepackage{newunicodechar}
\newunicodechar{→}{\ensuremath{\rightarrow}}
% Mismo problema con los simbolos de marca y vinetas: el «✓» de «una palomita ✓»
% salia como cuadrito vacio en el PDF de todas las guias que lo usaban.
\usepackage{pifont}
\newunicodechar{✓}{\ding{51}}
\newunicodechar{✗}{\ding{55}}
\newunicodechar{✔}{\ding{52}}
\newunicodechar{•}{\textbullet}
\newunicodechar{≥}{\ensuremath{\geq}}
\newunicodechar{≤}{\ensuremath{\leq}}

\setmainfont{Helvetica Neue}
\setsansfont{Helvetica Neue}
\setlength{\parindent}{0pt}
\setlength{\parskip}{6pt}
% Sin particion de palabras: en 6.º-7.º una palabra cortada por guion al final
% de linea («Comuni-cacion») frena la lectura mucho mas que una linea corta.
\hyphenpenalty=10000
\exhyphenpenalty=10000
\pretolerance=2000
\tolerance=3000
\setlength{\emergencystretch}{3em}
\linespread{1.10}
\renewcommand{\arraystretch}{1.25}
\setlist{leftmargin=5mm,itemsep=1mm,topsep=1mm}
\newcolumntype{Y}{>{\RaggedRight\arraybackslash}X}

\definecolor{milcMagenta}{HTML}{E6007E}
\definecolor{milcVino}{HTML}{5A0038}
\definecolor{milcTurquesa}{HTML}{0093A5}
\definecolor{milcVerde}{HTML}{58A923}
\definecolor{milcMostaza}{HTML}{E5B400}
\definecolor{milcOcre}{HTML}{B86600}
\definecolor{milcNegro}{HTML}{241816}
\definecolor{milcGris}{HTML}{F7F7F4}
\definecolor{milcLinea}{HTML}{E8E2DD}
\definecolor{milcGradePrimary}{HTML}{0D9488}
\definecolor{milcGradeDark}{HTML}{115E59}
\definecolor{milcGradeSoft}{HTML}{CCFBF1}
\definecolor{milcGradeAccent}{HTML}{CFE7FF}
\definecolor{milcGradeText}{HTML}{FFFFFF}
\color{milcNegro}

\pagestyle{fancy}
\fancyhf{}
% Cabecera de dos lineas: identidad de la guia arriba y, a la derecha y en
% gris, la estacion en curso (\markright desde \milcestacion). Antes de la
% primera estacion la segunda linea va vacia; el \strut mantiene la altura
% para que la linea de arriba no baile entre paginas.
\lhead{\small Territorio interior · Educación socioemocional\\[-1pt]{\footnotesize\strut}}
\rhead{\small Grado 11° · Momento 4 · MILC v3\\[-1pt]{\footnotesize\strut\color{milcNegro!65}\rightmark}}
\cfoot{\small\thepage}
\renewcommand{\headrulewidth}{0pt}

% Sobre mostaza (#E5B400) y verde (#58A923) el texto va en milcNegro; sobre el
% resto de la paleta, en blanco. Se usa dentro de fonttitle/fontupper, donde un
% \color posterior manda sobre coltitle.
\newcommand{\milctintasobre}[1]{%
  \ifboolexpr{test{\ifstrequal{#1}{milcMostaza}} or test{\ifstrequal{#1}{milcVerde}}}%
    {\color{milcNegro}}{\color{white}}}

\newtcolorbox{titlebox}[1]{
  enhanced,colback=#1,colframe=#1,arc=7mm,boxrule=0pt,
  left=7mm,right=7mm,top=3mm,bottom=3mm,
  before skip=8pt,after skip=8pt,
  fontupper=\sffamily\bfseries\Large\color{white}
}

\newtcolorbox{softbox}[3]{
  enhanced,breakable,pad at break=3mm,
  colback=#2,colframe=#1,borderline west={4pt}{0pt}{#1},
  arc=2mm,boxrule=.4pt,left=4mm,right=4mm,top=3mm,bottom=3mm,
  before skip=6pt,after skip=8pt,title={#3},coltitle=white,
  colbacktitle=#1,fonttitle=\sffamily\bfseries\milctintasobre{#1},fontupper=\RaggedRight,
  attach boxed title to top left={xshift=3mm,yshift=-2mm},
  boxed title style={arc=2mm, boxrule=0pt}
}

\newtcolorbox{drawbox}[2]{
  enhanced,breakable,pad at break=4mm,
  colback=white,colframe=#1,arc=4mm,boxrule=1pt,
  left=5mm,right=5mm,top=4mm,bottom=4mm,before skip=8pt,after skip=8pt,
  title={#2},coltitle=white,colbacktitle=#1,fonttitle=\sffamily\bfseries\milctintasobre{#1},
  fontupper=\RaggedRight,
  attach boxed title to top left={xshift=4mm,yshift=-2mm},
  boxed title style={arc=3mm, boxrule=0pt}
}

\newtcolorbox{infoband}[1]{
  enhanced,breakable,pad at break=2mm,colback=#1!8,colframe=#1!50,arc=2mm,boxrule=.5pt,
  left=4mm,right=4mm,top=2mm,bottom=2mm,before skip=4pt,after skip=4pt,
  fontupper=\small\RaggedRight
}

% Puente narrativo entre secciones (contrato editorial Conéctate).
% Da continuidad: "Acabas de X. Ahora vas a Y porque Z."
% Visualmente: banda izquierda en color del grado, texto en itálica pequeña.
\newtcolorbox{puentebox}{
  enhanced,breakable,pad at break=2mm,colback=milcGradeSoft,colframe=milcGradeDark,
  borderline west={3pt}{0pt}{milcGradeDark},
  arc=0mm,boxrule=0pt,left=5mm,right=4mm,top=2.5mm,bottom=2.5mm,
  before skip=4pt,after skip=8pt,
  fontupper=\itshape\small\color{milcNegro}\RaggedRight
}
% La flecha va en modo matematico a proposito: Helvetica Neue NO tiene el glifo
% U+2192, asi que \textrightarrow se compilaba a NADA («Missing character: There
% is no → in font Helvetica Neue») y el puente salia sin flecha. En modo
% matematico la flecha viene de la fuente matematica y si se dibuja.
\newcommand{\puente}[1]{%
  \begin{puentebox}%
    \textcolor{milcGradeDark}{$\rightarrow$}\hspace{2mm}#1%
  \end{puentebox}%
}

\newcommand{\linead}[1]{\noindent\color{milcLinea}\rule{#1}{0.5pt}\color{milcNegro}}
\newcommand{\respuesta}[1]{\par\vspace{2mm}\foreach \n in {1,...,#1}{\noindent\color{milcLinea}\rule{\linewidth}{0.5pt}\par\vspace{5mm}}\color{milcNegro}}
\newcommand{\checkbox}{\textcolor{milcGradeDark}{\fbox{\phantom{x}}}\hspace{5pt}}

% ─── Listas aireadas para las guías socioemocionales ────────────────────
% Componer las verificaciones como «\checkbox texto\\» deja el texto que salta
% de línea debajo del cuadrito y apelmaza el bloque. Estos entornos alinean la
% continuación y airean los ítems. Los usa Territorio Interior; cualquier
% familia puede hacerlo.
\newlist{checklista}{itemize}{1}
\setlist[checklista]{
  label=\textcolor{milcGradeDark}{\fbox{\phantom{x}}},
  leftmargin=9mm, labelsep=3.2mm, itemsep=2.4mm,
  topsep=2.5mm, parsep=0pt, partopsep=0pt, before=\RaggedRight
}
\newlist{pasoslista}{enumerate}{1}
\setlist[pasoslista]{
  label=\textcolor{milcGradeDark}{\sffamily\bfseries\arabic*.},
  leftmargin=8mm, labelsep=3mm, itemsep=2.8mm,
  topsep=1mm, parsep=0pt, partopsep=0pt, before=\RaggedRight
}

\newcommand{\phasecard}[4]{
\begin{tcolorbox}[enhanced,colback=#3,colframe=#2,arc=3mm,boxrule=.5pt,height=3.05cm,valign=center,halign=center,left=1.5mm,right=1.5mm,top=2mm,bottom=2mm]
{\sffamily\bfseries\fontsize{22}{24}\selectfont\textcolor{#2}{#1}}\par
{\sffamily\bfseries\small #4}
\end{tcolorbox}
}

% ── Piezas del rediseno 2026 (legibilidad 6.º-11.º) ───────────────────
% Banda de estacion: reemplaza el titlebox macizo. Titulo a la izquierda,
% dato practico (tiempo/modalidad) a la derecha, en una sola linea.
% Nunca huerfana: pide 9 lineas libres (o salta de pagina rellenando con
% \vfill) y prohibe el salto justo despues de la banda. Nueve y no seis:
% la caja partible que sigue necesita ~80pt (titulo adosado, relleno y dos
% lineas) para arrancar en la misma pagina; con menos, tcolorbox la manda
% entera a la siguiente y la banda se queda sola. El penalty 10000 va
% en `after`, ANTES del espacio posterior: un `after skip` (aunque sea 0pt)
% es un \vskip tras la caja, y ese glue es un punto de corte legal que un
% \nopagebreak escrito despues de \end{tcolorbox} ya no cierra. Cada banda
% deja marcador PDF y actualiza la cabecera derecha.
\newcounter{milcestacion}
\newcommand{\milcestacion}[2]{%
\Needspace*{9\baselineskip}%
\stepcounter{milcestacion}%
\pdfbookmark[1]{#2}{milcestacion\themilcestacion}%
\markright{#2}%
\begin{tcolorbox}[enhanced,colback=#1,colframe=#1,arc=2mm,boxrule=0pt,
  left=5mm,right=5mm,top=2.6mm,bottom=2.6mm,before skip=11pt,
  after={\par\nopagebreak\vskip7pt}]
{\sffamily\bfseries\fontsize{15}{18}\selectfont\milctintasobre{#1}#2}
\end{tcolorbox}}

% Tarjeta de paso numerada: sustituye las tablas «Pilar / Aplicacion», que
% eran formato de consulta docente, no de estudiante. El numero va en un
% circulo sobre el borde para que la secuencia se lea de un vistazo.
\newcommand{\milcpaso}[3]{%
\Needspace{4\baselineskip}%
\begin{tcolorbox}[enhanced,colback=white,colframe=milcLinea,arc=1.5mm,boxrule=.9pt,
  left=14mm,right=4mm,top=3mm,bottom=3mm,before skip=4pt,after skip=4pt,
  fontupper=\RaggedRight,
  overlay={\node[anchor=north west,circle,fill=#2,text=white,inner sep=0pt,
    minimum size=7.4mm,font=\sffamily\bfseries\small]
    at ([xshift=3.6mm,yshift=-3.2mm]frame.north west) {#1};}]
#3
\end{tcolorbox}}

% Casilla marcable de tamano util con lapiz (la anterior era \fbox{\phantom{x}},
% de alto variable segun el texto que la seguia).
\newcommand{\milcmarca}{\framebox[3.8mm]{\rule{0pt}{3.4mm}}}

% Fila marcable de lista suelta (checks de la estacion 1).
\newcommand{\milccheckrow}[1]{%
\par\vspace{1.8mm}\noindent
\makebox[6.4mm][l]{\raisebox{-.55mm}{\textcolor{milcNegro}{\milcmarca}}}%
\parbox[t]{\dimexpr\linewidth-6.4mm\relax}{\RaggedRight #1}\par}

% Celda marcable dentro de la rubrica (tabla de autoevaluacion). Misma
% sangria francesa, pero pensada para vivir dentro de una columna X.
\newcommand{\milcopcion}[1]{%
\makebox[5.2mm][l]{\raisebox{-.55mm}{\textcolor{milcNegro}{\milcmarca}}}%
\parbox[t]{\dimexpr\linewidth-5.2mm\relax}{\RaggedRight #1}}

% ── Recursos visuales de la guía ──────────────────────────────────────
% Declarados en el bloque `recursos:` del YAML y emitidos por el builder.
% \guiaFigura   → imagen rasterizada o PDF (foto, carta celeste, montaje).
% \guiaDiagrama → fuente TikZ/circuitikz (.tex) incrustada como vector:
%                 es la vía para circuitos y esquemáticos editables.
% [#1] = texto alternativo (alt) · #2 = ruta relativa al .tex · #3 = pie
% El alt llega al PDF etiquetado (/Alt) y lo lee el lector de pantalla.
\newcommand{\guiaFigura}[3][]{%
\begin{center}
\includegraphics[width=0.88\linewidth,height=0.38\textheight,keepaspectratio,alt={#1}]{#2}\\[1.5mm]
{\footnotesize\itshape\textcolor{milcNegro!75}{#3}}
\end{center}
}

\newcommand{\guiaDiagrama}[2]{%
\begin{center}
\input{#1}\\[1.5mm]
{\footnotesize\itshape\textcolor{milcNegro!75}{#2}}
\end{center}
}

\begin{document}
\thispagestyle{empty}

% ── Portada ───────────────────────────────────────────────────────────
% Antes: un rectangulo de color a pagina completa. Se veia bien en pantalla,
% pero en la fotocopiadora del colegio se comia el toner y salia gris sucio.
% Ahora la banda ocupa el tercio superior y el resto va en blanco.
\begin{tikzpicture}[remember picture,overlay]
  \path[fill=milcGradePrimary]
    (current page.north west) rectangle ([yshift=-10.6cm]current page.north east);

  \node[anchor=north west,text=milcGradeText] at ([xshift=2.0cm,yshift=-1.55cm]current page.north west)
    {\sffamily\bfseries\fontsize{12}{14}\selectfont MOMENTO};
  \node[anchor=north west,text=milcGradeText,opacity=.85] at ([xshift=2.0cm,yshift=-2.35cm]current page.north west)
    {\sffamily\bfseries\fontsize{8.5}{10}\selectfont TERRITORIO INTERIOR · SOCIOEMOCIONAL};
  \node[anchor=north east,text=milcGradeText,opacity=.85,align=right] at ([xshift=-2.0cm,yshift=-1.55cm]current page.north east)
    {\sffamily\bfseries\fontsize{9}{12}\selectfont I.E. Sor María Juliana\\Cartago, Valle del Cauca};

  \node[anchor=west,text=milcGradeText] (gradonum) at ([xshift=1.85cm,yshift=-7.1cm]current page.north west)
    {\sffamily\bfseries\fontsize{112}{112}\selectfont 4};
  % El circulito de grado (°) solo aplica a las guias de grado («11°»); en
  % Bebras y Semillero el numero grande es un momento/modulo, no un grado.
  % Se posiciona relativo al nodo `gradonum`, no en coordenadas absolutas.
  

  \node[anchor=west,text=milcGradeText,text width=11.4cm,align=left] at ([xshift=1.5cm,yshift=0.5cm]gradonum.east)
    {
     {\sffamily\bfseries\fontsize{9}{11}\selectfont\MakeUppercase{Territorio interior · Socioemocional}}\\[3mm]
     {\sffamily\bfseries\fontsize{21}{25}\selectfont\setlength{\baselineskip}{27pt}El trabajo\\[4pt]y la dignidad\\[4pt]tres maneras de vivir un oficio}\\[3.5mm]
     {\sffamily\bfseries\fontsize{15}{18}\selectfont Grado 11° · Momento 4}
    };
\end{tikzpicture}

\vspace*{9.4cm}
\vfill

{\sffamily\bfseries\fontsize{8}{10}\selectfont\textcolor{milcGradeDark}{LO QUE TE LLEVAS HOY}}

\begin{tcolorbox}[enhanced,colback=white,colframe=milcNegro,arc=2mm,boxrule=1.6pt,
  left=5mm,right=5mm,top=4mm,bottom=4mm,before skip=3pt,after skip=12pt,
  fontupper=\RaggedRight\fontsize{11}{15.5}\selectfont]
Tu mapa del oficio: los tres trabajos que conoces de cerca clasificados en trabajo, carrera o vocación con la evidencia de cada uno, las cuatro preguntas aplicadas a lo que quieres hacer, y una conversación real con alguien que ya vive de eso.
\end{tcolorbox}

\begin{tcolorbox}[enhanced,colback=milcGris,colframe=milcGris,arc=2mm,boxrule=0pt,
  left=5mm,right=5mm,top=3.5mm,bottom=3.5mm,after skip=12pt]
{\sffamily\bfseries\fontsize{8}{10}\selectfont\textcolor{milcNegro!55}{ESTA GUÍA ES DE}}\\[3mm]
\begin{tabularx}{\linewidth}{X p{3.2cm} p{3.2cm}}
\linead{\linewidth} & \linead{3.0cm} & \linead{3.0cm}\\[-1mm]
{\fontsize{8.5}{10}\selectfont\textcolor{milcNegro!55}{Nombre completo}} &
{\fontsize{8.5}{10}\selectfont\textcolor{milcNegro!55}{Grupo}} &
{\fontsize{8.5}{10}\selectfont\textcolor{milcNegro!55}{Fecha}}\\
\end{tabularx}
\end{tcolorbox}

\vfill

\noindent\textcolor{milcLinea}{\rule{\linewidth}{1pt}}\\[2mm]
{\sffamily\bfseries\fontsize{9.5}{12}\selectfont Autor: Dr. Álvaro Cárdenas Orozco}\\[1mm]
{\sffamily\fontsize{8.5}{11}\selectfont\textcolor{milcNegro!55}{Metodología MILC · Apertura ancestral · Trabajo y dignidad · Triángulo Dussel-Estoicismo-Floridi}}

\newpage

\section*{Apertura: El trabajo y la dignidad: el mismo oficio, tres maneras de vivirlo}

\begin{softbox}{milcGradeDark}{milcGradeSoft}{Saber ancestral}
En el Norte del Valle, una pieza de calado \textbf{se reconoce}. \emph{Quien sabe mirar distingue la mano:} \textbf{la tensión del hilo, el remate, la manera de resolver una esquina, la puntada preferida.} \emph{Y no es una metáfora ---en los talleres se sabe de qué mano salió una pieza, y se sabe a quién encargarle qué---.} \textbf{Ese es el rasgo de lo que se llama \emph{un oficio con nombre propio}:} \emph{el trabajo no es anónimo, sale con la marca de quien lo hizo, y por eso la persona responde por él.} \textbf{Fíjate en las tres cosas que eso produce, porque las tres tienen que ver con la dignidad.} \textbf{Primera: hay criterio propio} ---la caladora sabe cuándo su pieza está bien hecha sin que nadie se lo diga, y ya viste en el momento 1 que eso se juega en el revés---. \textbf{Segunda: hay mejora posible} ---se puede llegar a hacerlo mejor que ayer, y eso se nota---. \textbf{Tercera: hay alguien del otro lado} ---la pieza va a usarse, y quien la hizo lo sabe---. \emph{Y aquí conviene decir algo sin romanticismo:} \textbf{el calado es un trabajo mal pagado, y esta guía no va a fingir que no lo es.} \emph{Precisamente por eso sirve de ancla: muestra que la dignidad de un oficio y lo que ese oficio paga son dos cosas distintas ---y que confundirlas hace daño en las dos direcciones---.}
\end{softbox}

\begin{softbox}{milcTurquesa}{milcGris}{Saber contemporáneo}
¿Y de qué depende que alguien la pase bien en su trabajo? \textbf{Menos de la ocupación de lo que uno creería.} \textbf{Amy Wrzesniewski} y su equipo estudiaron a \textbf{196 empleados} de dos lugares de trabajo, con ocupaciones que iban de lo administrativo a lo profesional, y encontraron que la gente ve su trabajo \textbf{de una de tres maneras, sin ambigüedad}. Como un \textbf{empleo} ---el foco está en la retribución económica y en la necesidad, no en el disfrute; no es una parte importante de la vida---. Como una \textbf{carrera} ---el foco está en el avance, en subir---. O como una \textbf{vocación} ---el foco está en el disfrute de un trabajo que llena y que es socialmente útil--- (Wrzesniewski, McCauley, Rozin y Schwartz, 1997). \textbf{Y aquí está el hallazgo que conviene subrayar, porque desarma la manera habitual de elegir carrera:} \emph{esas diferencias \textbf{no se reducían a la ocupación ni a los datos demográficos}}. \textbf{En un subgrupo homogéneo de \textbf{24 asistentes administrativas} ---el mismo cargo, el mismo lugar--- aparecían igual las tres orientaciones.} \emph{Es decir: el mismo trabajo puede vivirse como empleo, como carrera o como vocación.} \textbf{No lo decide el oficio: lo decide la relación con él.}
\end{softbox}

\begin{softbox}{milcMostaza}{milcGris}{Pregunta-puente}
\textit{Una pieza de calado \textbf{se reconoce}, y su oficio está mal pagado. \textbf{¿Qué trabajo estás evaluando por lo que paga o por lo que dirían} \ldots \textbf{sin haber visto nunca cómo es un día entero de esa vida}?}
\end{softbox}

\begin{softbox}{milcVerde}{milcGris}{Saber y saber hacer}
Al terminar podrás: (1) \textbf{distinguir} las tres relaciones con el trabajo ---empleo, carrera, vocación--- y reconocer que el mismo oficio admite las tres; (2) \textbf{entender} que la dignidad de un trabajo y su remuneración son cosas distintas, sin usar eso para justificar que se pague mal; (3) \textbf{evaluar} un oficio por cómo es un día completo y no por su nombre; (4) \textbf{hablar} con alguien que ya vive de lo que te interesa, y traer lo que averiguaste.
\end{softbox}



\small
\begin{tabularx}{\textwidth}{>{\bfseries}p{3.0cm}Y}
\rowcolor{milcGradePrimary!12}Autor & Dr. Álvaro Cárdenas Orozco\\
DBA & Comprende los fenómenos que afectan a su territorio, regula sus emociones ante ellos y acompaña a otros con empatía (transversal 6°--11°, articulado con la política «Escuela, territorio de vida» del MEN, Resolución 006519 de 2025, y la Ley 1620 de 2013)\\
Referentes & Primera ayuda psicológica (OMS, 2011/2012) · Directrices IASC (2007) · USGS y Servicio Geológico Colombiano · Pueblo nasa y la minga (CRIC) · Dussel · Estoicismo · Floridi\\
Producto final & Tu mapa del oficio: los tres trabajos que conoces de cerca clasificados en trabajo, carrera o vocación con la evidencia de cada uno, las cuatro preguntas aplicadas a lo que quieres hacer, y una conversación real con alguien que ya vive de eso.\\
\end{tabularx}
\normalsize

\puente{En el momento 3 trabajaste cómo decidir. \textbf{Hoy miras la decisión que más pesa a los diecisiete}, \emph{y la miras por donde casi nadie la mira: no por el nombre del oficio, sino por cómo se vive por dentro}.}

\milcestacion{milcGradeDark}{Ruta MILC: así vamos a aprender}
\begin{tabularx}{\textwidth}{>{\centering\arraybackslash}X>{\centering\arraybackslash}X>{\centering\arraybackslash}X>{\centering\arraybackslash}X}
\phasecard{1}{milcVerde}{milcGris}{Escucha\\[-1mm]\small Observo, escucho y pregunto.} &
\phasecard{2}{milcTurquesa}{milcGris}{Sistematización\\[-1mm]\small Distingo el oficio de cómo se vive.} &
\phasecard{3}{milcMagenta}{milcGris}{Praxis\\[-1mm]\small Construyo y aplico.} &
\phasecard{4}{milcVino}{milcGris}{Evaluación\\[-1mm]\small Reflexión liberadora.}
\end{tabularx}


\puente{Primero, los trabajos que conoces de cerca. \emph{No los que suenan bien: los que has visto de verdad.}}

\milcestacion{milcVerde}{Estación 1 de 3 · Escucha}

\begin{softbox}{milcVerde}{milcGris}{Escena para escuchar}
\textbf{Actividad 1 · IDENTIFICA --- Los trabajos que he visto.} \textbf{Sin nombres: usa iniciales.} \textbf{Primera parte.} Escribe \textbf{tres trabajos que conozcas de cerca} ---de tu casa, de tu barrio, de alguien que quieres---. \emph{No los que suenan bien: los que de verdad has visto de cerca.} \textbf{Segunda parte.} De cada uno responde: \emph{¿esa persona la pasa bien? ¿cómo lo sabes? ¿qué es lo peor de ese trabajo? ¿qué es lo mejor?} \textbf{Y una que casi nadie hace:} \emph{¿esa persona sabe cuándo hizo bien su trabajo?} \textbf{Tercera parte.} Escribe \textbf{qué te gustaría hacer} ---o dos opciones, si no lo tienes claro--- y responde con honestidad: \emph{¿por qué eso?} \textbf{Marca la razón real:} \emph{me gusta · paga bien · es lo que esperan de mí · es lo que puedo · no sé.} \textbf{Y la última:} \emph{¿has visto alguna vez un día completo de trabajo de alguien que haga eso?}
\end{softbox}

{\sffamily\bfseries\fontsize{8}{10}\selectfont\textcolor{milcNegro!55}{MARCA CUANDO LO HAYAS HECHO}}
\milccheckrow{Escogiste tres trabajos que has visto de cerca, no que suenan bien.}
\milccheckrow{Respondiste si esa persona \textbf{sabe cuándo hizo bien su trabajo}.}
\milccheckrow{Marcaste la razón real por la que te interesa lo que te interesa.}

\begin{infoband}{milcVerde}


\textbf{En tu cuaderno (Actividad 1 · IDENTIFICA).} Título: «Actividad 1. Los trabajos que he visto». \textbf{Formato:} los tres trabajos con sus respuestas + lo que te gustaría hacer y la razón real + si has visto un día completo. \textbf{Extensión:} media página. \textbf{Sabes que terminaste cuando} marcaste \textbf{la razón real}. Esta página es tuya: \textbf{nadie más tiene que leerla si tú no quieres}.
\end{infoband}

\puente{Ya los tienes. Ahora, las tres maneras de vivir un trabajo y por qué no las decide la ocupación.}

\milcestacion{milcTurquesa}{Fase 2 · Sistematización: entender el oficio}

\begin{softbox}{milcTurquesa}{milcGris}{Cuatro claves sobre el trabajo}
Cuatro claves para evaluar un oficio por algo distinto de su nombre.
\end{softbox}

\milcpaso{1}{milcTurquesa}{\textbf{Tres relaciones, no tres tipos de empleo.} \textbf{Empleo:} el foco está en la plata y en la necesidad; el trabajo no es una parte importante de la vida. \textbf{Carrera:} el foco está en el avance, en llegar más arriba. \textbf{Vocación:} el foco está en el disfrute de algo que llena y que es \textbf{socialmente útil} (Wrzesniewski y otros, 1997). \emph{Y conviene decirlo sin moralina:} \textbf{ninguna de las tres es despreciable}. \emph{Quien trabaja para comer no está haciendo algo menor ---está sosteniendo una casa---,} \textbf{y buena parte del mundo funciona porque hay gente que va a trabajar sin que le encante.} \emph{Lo que sí conviene saber es \textbf{cuál es la tuya}, porque cada una tiene un costo distinto:} \textbf{el empleo cobra en tiempo, la carrera cobra en comparación} ---ya viste el momento 2--- \emph{y la vocación cobra porque cuesta poner límites cuando uno ama lo que hace.}}
\milcpaso{2}{milcTurquesa}{\textbf{No lo decide la ocupación, y esto cambia cómo elegir.} \emph{El dato más contundente del estudio:} \textbf{las diferencias no se reducían a la ocupación ni a lo demográfico}, y \emph{en un subgrupo homogéneo de \textbf{24 asistentes administrativas} ---mismo cargo, mismo sitio--- aparecían las tres orientaciones}. \textbf{El mismo trabajo, vivido de tres maneras.} \emph{Piensa lo que eso implica para la pregunta «¿qué carrera me hará feliz?»:} \textbf{está mal formulada}. \emph{No hay carreras que garanticen una relación buena con el trabajo, ni carreras que la impidan.} \textbf{Hay condiciones que ayudan} ---que uno vea el resultado, que pueda mejorar, que sepa a quién le sirve--- \emph{y esas condiciones se dan o no se dan \textbf{dentro} de casi cualquier oficio.}}
\milcpaso{3}{milcTurquesa}{\textbf{La firma de la pieza: tres condiciones que sí importan.} \emph{Vuelve al calado y mira qué tiene ese oficio, aunque pague poco.} \textbf{Criterio propio:} la caladora \emph{sabe} cuándo su pieza quedó bien, sin que nadie se lo diga. \textbf{Mejora visible:} se puede hacerlo mejor que el año pasado, y se nota. \textbf{Alguien del otro lado:} la pieza va a usarse, y quien la hizo lo sabe. \emph{Esas tres condiciones ---saber cuándo lo hiciste bien, poder mejorar, saber a quién le sirve--- son las que separan un trabajo que sostiene de uno que vacía,} \textbf{y son mucho mejores preguntas que «¿cuánto paga?» a la hora de mirar un oficio.} \emph{Un trabajo bien pagado donde nunca sabes si lo hiciste bien es más difícil de sostener de lo que parece a los diecisiete.}}
\milcpaso{4}{milcTurquesa}{\textbf{La dignidad no es el salario ---y decirlo no justifica que se pague mal---.} \emph{Hay que sostener las dos mitades de esta frase, porque la gente suele soltar una.} \textbf{Primera mitad:} \emph{no hay oficios indignos}; \textbf{quien barre, quien cuida, quien cocina, quien recoge café, quien borda, hace un trabajo del que depende que otros vivan}, y despreciarlo es la misma lógica de clase del momento 5 del grado 9. \textbf{Segunda mitad, que casi siempre se omite:} \emph{que un oficio sea digno \textbf{no es una razón para que pague mal}}. \textbf{El calado paga poco, y eso no es una virtud del oficio: es una injusticia del mercado que lo compra.} \emph{Confundir las dos cosas ---«es que ellas lo hacen por amor al arte»--- es el argumento con que históricamente se le ha pagado menos al trabajo de las mujeres.} \textbf{Dignidad y remuneración son cosas distintas, y las dos se reclaman.}}

\begin{drawbox}{milcTurquesa}{Las cuatro preguntas para mirar un oficio}
\begin{pasoslista}
\item \textbf{¿Sabría cuándo lo hice bien?} \emph{¿Hay criterio propio, o dependo de que alguien me califique?}
\item \textbf{¿Puedo mejorar, y se nota?} \emph{Un trabajo donde el año diez es igual al año uno cansa distinto.}
\item \textbf{¿A quién le sirve lo que hago?} \emph{Poder nombrar a alguien concreto cambia el día.}
\item \textbf{¿Cómo es un martes cualquiera?} \emph{No el día bueno ni el día malo: el corriente.}
\item \textbf{Y la de realidad, que no se omite:} \emph{¿de qué se vive haciendo esto, y alcanza?}
\end{pasoslista}
\end{drawbox}

\begin{softbox}{milcMostaza}{milcGris}{Errores comunes y su corrección}
\textbf{Evaluar por el nombre del oficio:} el mismo cargo admite las tres relaciones. \textbf{«¿Qué carrera me hará feliz?»:} la pregunta está mal formulada. \textbf{Despreciar oficios manuales o de cuidado:} de ellos depende que otros vivan. \textbf{Usar «es digno» para justificar que pague mal:} es el argumento con que se ha pagado menos a las mujeres. \textbf{Mirar solo el día bueno:} el oficio es el martes corriente. \textbf{Preguntarle a quien no lo ejerce:} los folletos no cuentan la rutina. \textbf{Elegir solo por lo que pagan:} el criterio propio y la mejora sostienen más de lo que parece.
\end{softbox}

\begin{infoband}{milcTurquesa}


\textbf{En tu cuaderno (Actividad 2 · EXPLICA).} Título: «Actividad 2. Empleo, carrera, vocación». \textbf{Formato:} define las tres con tus palabras y clasifica los tres trabajos de la Actividad 1, \textbf{con la evidencia} en que te basas ---qué dijo o hizo esa persona---. \textbf{Extensión:} media página. \textbf{Sabes que terminaste cuando} puedes explicar por qué \textbf{24 personas con el mismo cargo tenían las tres orientaciones}.
\end{infoband}

\puente{Ya sabes qué mirar. Es hora de mirar un día completo y de hablar con alguien que ya está adentro.}

\milcestacion{milcMagenta}{Fase 3 · Praxis: mirar de cerca}

\begin{softbox}{milcMagenta}{milcGris}{Cómo se averigua de verdad qué es un oficio}
Ningún folleto cuenta cómo es un martes. Esta actividad se hace preguntándole a alguien que ya está adentro.
\end{softbox}

\milcpaso{1}{milcMagenta}{\textbf{Los tres trabajos, clasificados (7 min).} Pon los tres en las tres categorías \textbf{con evidencia observable} ---«dice que cuenta los días para el viernes» es evidencia; «se le nota» no lo es---. \emph{Y fíjate si alguno te sorprende:} \textbf{es común descubrir que alguien que uno creía resignado tiene una relación de vocación con su oficio, y al revés.}}
\milcpaso{2}{milcMagenta}{\textbf{Las cuatro preguntas (8 min).} Aplícalas a \textbf{lo que te gustaría hacer}. \emph{Vas a encontrar que a la mayoría no puedes responderlas} ---y eso no es un fracaso: \textbf{es exactamente el descubrimiento que hace falta}---. \emph{Marca cuáles no sabes: esas son las que vas a preguntar.}}
\milcpaso{3}{milcMagenta}{\textbf{El día completo (7 min).} Escribe \textbf{cómo te imaginas un martes corriente} de alguien que hace lo que te interesa: \emph{a qué hora empieza, qué hace las primeras dos horas, con quién habla, qué es lo más aburrido, qué es lo más difícil, a qué hora termina.} \emph{Es la técnica del momento 1 del grado 9 aplicada a un oficio,} \textbf{y como allá, lo que enseña es en qué te equivocaste.}}
\milcpaso{4}{milcMagenta}{\textbf{La conversación (fuera de clase).} \textbf{Habla con alguien que ya viva de eso} ---no un profesor de la materia: alguien que trabaje en eso---. \emph{Si no conoces a nadie, pregunta en tu casa, en el colegio o escribe:} \textbf{la mayoría de la gente contesta cuando alguien de diecisiete le pregunta en serio por su oficio.} \textbf{Cuatro preguntas y nada más:} \emph{«¿cómo es un martes cualquiera?», «¿cómo sabe usted cuándo hizo bien su trabajo?», «¿qué es lo que menos le gusta?», «¿de qué se vive haciendo esto?»}. \textbf{Y después compara con tu día imaginado} y anota \textbf{en qué te equivocaste}.}

\begin{drawbox}{milcMagenta}{Antes de cerrar tu mapa del oficio}
\begin{checklista}
\item Clasificaste los tres trabajos \textbf{con evidencia observable}.
\item Aplicaste las cuatro preguntas y \textbf{marcaste las que no sabes responder}.
\item Escribiste tu martes imaginado, con detalle y sin adornar.
\item \textbf{Hablaste con alguien que vive de eso}, no con alguien que lo enseña.
\item Anotaste \textbf{en qué te equivocaste} al imaginar el día.
\item Incluiste la pregunta de realidad: de qué se vive haciendo eso.
\end{checklista}
\end{drawbox}

\begin{softbox}{milcMagenta}{milcGris}{Plantilla de trabajo}
\textbf{Las cuatro preguntas.} \emph{«¿Sabría cuándo lo hice bien? · ¿Puedo mejorar y se nota? · ¿A quién le sirve? · ¿Cómo es un martes?»} \\[1mm]
\textbf{Para la conversación.} \emph{«¿Cómo es un martes cualquiera? · ¿Cómo sabe cuándo hizo bien su trabajo? · ¿Qué es lo que menos le gusta? · ¿De qué se vive haciendo esto?»} \\[1mm]
\textbf{El registro.} \emph{«Creí que \_\_\_\_. En realidad \_\_\_\_.»}
\end{softbox}

\begin{infoband}{milcMagenta}


\textbf{En tu cuaderno (Actividad 3 · CREA).} Título: «Actividad 3. Mi mapa del oficio». \textbf{Formato:} los tres trabajos clasificados con evidencia + las cuatro preguntas aplicadas + el martes imaginado + lo que respondió la persona y en qué te equivocaste. \textbf{Extensión:} una página. \textbf{Sabes que terminaste cuando} \textbf{ya tuviste la conversación}.
\end{infoband}

\puente{El producto es una conversación real. \emph{Ningún folleto de universidad cuenta cómo es un martes cualquiera en ese oficio.}}

\milcestacion{milcMostaza}{Producto de la sesión}
\textbf{Mapa del oficio: tres trabajos clasificados con evidencia, las cuatro preguntas y una conversación con quien ya vive de eso}.
\begin{softbox}{milcMostaza}{milcGris}{Criterios de calidad}
(1) La clasificación en empleo, carrera y vocación se apoya en evidencia observable, no en impresiones. (2) Las cuatro preguntas están aplicadas y se identifican explícitamente las que no se pueden responder aún. (3) El día imaginado es concreto y describe una jornada corriente, no la excepcional. (4) La conversación se hizo con alguien que ejerce el oficio, no con quien lo enseña. (5) Se registran las diferencias entre lo imaginado y lo que contó esa persona.
\end{softbox}

\puente{Antes de cerrar, revisa lo importante: si tu evaluación de un oficio se basa en lo que paga y en lo que dirían, no averiguaste nada nuevo.}

\milcestacion{milcVino}{Evaluación liberadora · cinco dimensiones}

\begin{softbox}{milcVino}{milcGris}{Sentido de la evaluación}
Evaluar no es solo revisar una nota. Es reconocer qué aprendí, qué cambió en mí, qué emociones manejé, cómo me relaciono con la ciudad y la comunidad, y qué le dejaré a quien viene detrás.
\end{softbox}

\textbf{Mi reflexión liberadora en cinco dimensiones.} Completa cada frase iniciada:

\small
\begin{tabularx}{\textwidth}{>{\bfseries\RaggedRight\arraybackslash}p{3.4cm}Y>{\RaggedRight\arraybackslash}p{4.6cm}}
\rowcolor{milcVino}\color{white}Dimensión & \color{white}\bfseries Mi respuesta (frame starter) & \color{white}\bfseries Algo que descubrí\\
1. Personal & Lo que más cambió en mí fue \linead{6.5cm} & \linead{4.2cm}\\
2. Emocional & La emoción más fuerte fue \linead{2.5cm} y la regulé \linead{3cm} & \linead{4.2cm}\\
3. Ciudadana & En democracia esto importa porque \linead{6cm} & \linead{4.2cm}\\
4. Local (Cartago) & En mi barrio sería útil para \linead{6.5cm} & \linead{4.2cm}\\
5. Intergeneracional & A mi abuela/abuelo le diría \linead{6.5cm} & \linead{4.2cm}\\
\end{tabularx}
\normalsize

\begin{infoband}{milcVino}
\textbf{Para la próxima sustentación.} Vuelve a esta tabla en seis meses. Verás que las respuestas cambian — no porque lo aprendido cambie, sino porque tú cambias. La evaluación liberadora es una conversación contigo a través del tiempo.
\end{infoband}


\puente{Cerramos con tres voces: el rostro que reclama un derecho, la armonía entre palabras y acciones, y el trabajo que hoy se hace con información.}

\milcestacion{milcGradeDark}{Triángulo de pensamiento · cierre filosófico}

\begin{softbox}{milcVino}{milcGris}{Dussel · lente del nosotros}
\textit{«El otro se revela realmente como otro\ldots como el pobre, el oprimido; el que, a la vera del camino, fuera del sistema, muestra su rostro sufriente y sin embargo desafiante: «¡Tengo hambre!, ¡tengo derecho a comer!»»}

{\footnotesize\itshape\color{milcNegro!70}--- Enrique Dussel · Filosofía de la liberación (1977)}

\emph{Esta cita sostiene la segunda mitad del cuarto pilar, la que casi siempre se omite.} \textbf{Decir que un oficio es digno no puede servir para justificar que pague mal.} \emph{Las caladoras de Cartago hacen un trabajo reconocible, técnico y patrimonial ---y viven de él con dificultad---.} \textbf{Y la última palabra de la cita vuelve a ser la clave:} \emph{lo que se reclama no es admiración por el oficio, es un derecho}. \textbf{Fíjate en el argumento que se usa para pagar poco:} \emph{«es que lo hacen por amor», «es que es su pasión».} \textbf{Ese argumento se ha usado históricamente, y sobre todo con el trabajo de las mujeres y con el trabajo de cuidado.} \emph{Reconocer la dignidad de un oficio y exigir que se pague bien no son cosas opuestas: son la misma exigencia.}

\textbf{¿Qué trabajos de los que sostienen mi vida diaria están peor pagados de lo que deberían?}
\end{softbox}
\linead{\linewidth}\\[1mm]
\linead{\linewidth}

\begin{softbox}{milcMostaza}{milcGris}{Estoicismo · Séneca · Cartas a Lucilio, 20 (c. 64 d.C.) · cuidado interior}
\textit{«La filosofía enseña a obrar, no a hablar; quiere que cada uno viva a la manera que prescribe, que estén en armonía nuestras palabras con nuestras acciones.»}

{\footnotesize\itshape\color{milcNegro!70}--- Séneca · Cartas a Lucilio, 20 (c. 64 d.C.)}

\emph{Séneca describe, sin saberlo, lo que Wrzesniewski llamó \textbf{vocación}:} \textbf{que lo que uno hace y lo que uno es no vayan por caminos distintos.} \emph{Y conviene aterrizarlo, porque «vocación» suena a algo místico que a uno le llega:} \textbf{en el estudio, la vocación no era una ocupación ---era una relación---}, \emph{y aparecía en cargos que nadie llamaría vocacionales.} \textbf{Es decir: no se encuentra, se construye}, \emph{y se construye con lo que dice el pilar tres ---criterio propio, mejora, alguien del otro lado---.} \emph{Y hay un aviso honesto:} \textbf{el que ama lo que hace tiene la dificultad de poner límites}, \emph{y eso ya lo trabajaste en el momento 10 del grado 10.}

\textbf{Lo que quiero hacer, ¿lo quiero por lo que es o por lo que significaría ante otros?}
\end{softbox}
\linead{\linewidth}\\[1mm]
\linead{\linewidth}

\begin{softbox}{milcTurquesa}{milcGris}{Floridi · lente de la infoesfera}
\textit{«Somos organismos informacionales ---inforgs--- que compartimos un entorno global hecho, en última instancia, de información: la infoesfera.»}

{\footnotesize\itshape\color{milcNegro!70}--- Luciano Floridi · La cuarta revolución (2014)}

\textbf{Buena parte de los oficios que van a existir cuando ustedes salgan todavía no tiene nombre, y varios de los que hoy suenan seguros van a cambiar mucho.} \emph{Eso asusta, y hay una manera razonable de mirarlo:} \textbf{si el nombre del oficio es lo que cambia, entonces apostarle al nombre es lo frágil.} \emph{Lo que se mueve mejor entre oficios es otra cosa:} \textbf{saber aprender, saber trabajar con otros, saber explicar lo que uno hace, saber preguntar bien.} \emph{Y una observación específica sobre el trabajo de hoy:} \textbf{cuando el resultado es información y no una pieza, es más difícil saber cuándo uno hizo bien su trabajo} ---no hay un revés que revisar---. \emph{Por eso, en un oficio de pantalla, hay que \textbf{construirse} el criterio propio a propósito.}

\textbf{Si el oficio que quiero cambia de nombre en diez años, ¿qué de lo que aprendí seguiría sirviéndome?}
\end{softbox}
\linead{\linewidth}\\[1mm]
\linead{\linewidth}

\begin{infoband}{milcNegro}
\textbf{Nota para el docente.} Las tres son citas textuales y llevan su referencia debajo. Puedes remitir a la obra en clase.
\end{infoband}

\milcestacion{milcVino}{Mi autoevaluación · marca una casilla por fila}

{\small
\setlength{\extrarowheight}{2.2pt}
\begin{tabularx}{\textwidth}{>{\RaggedRight\arraybackslash\bfseries}p{3.0cm}YYY}
\rowcolor{milcVino}\color{white}\bfseries Criterio & \color{white}\bfseries Lo logré & \color{white}\bfseries En proceso & \color{white}\bfseries Necesito apoyo\\[.6mm]
Comprensión del tema & \milcopcion{Explico las ideas con mis palabras.} & \milcopcion{Reconozco las ideas, falta precisión.} & \milcopcion{Necesito repasar lo básico.}\\[.8mm]
\rowcolor{milcGris}Calidad de la indagación & \milcopcion{Distingo la ocupación de la relación con ella, y hablé con alguien que vive de lo que me interesa.} & \milcopcion{Entiendo la diferencia, pero sigo evaluando oficios por lo que pagan o por su prestigio.} & \milcopcion{Sigo pensando que hay trabajos dignos y trabajos que no lo son.}\\[.8mm]
Aplicación y producto & \milcopcion{Mi producto cumple los criterios.} & \milcopcion{Está armado, con ajustes pendientes.} & \milcopcion{Aún está en construcción.}\\[.8mm]
\rowcolor{milcGris}Reflexión 5 dimensiones & \milcopcion{Respondí cada una con honestidad.} & \milcopcion{Respondí algunas con detalle.} & \milcopcion{Mis respuestas son muy generales.}\\[.8mm]
Triángulo de pensamiento & \milcopcion{Conecté las 3 voces con mi trabajo.} & \milcopcion{Conecté al menos una.} & \milcopcion{No las relaciono con mi proceso.}\\[.8mm]
\end{tabularx}}

\vspace{3mm}
\textbf{Hoy entendí que:}\\[1mm]
\linead{\linewidth}\\[1mm]
\linead{\linewidth}

\textbf{Mi compromiso para la próxima sesión es:}\\[1mm]
\linead{\linewidth}\\[1mm]
\linead{\linewidth}

\begin{softbox}{milcMostaza}{milcGris}{Cita de despedida · Marco Aurelio}
\textit{``El obstáculo en el camino se vuelve el camino. Lo que estorba al camino se vuelve camino.''}\\[1mm]
Cada error de hoy, bien observado, se convierte en pieza del camino que pisarás mañana.
\end{softbox}

\small
\begin{tabularx}{\textwidth}{>{\bfseries}p{3.6cm}Y}
\rowcolor{milcGradeDark!12}Nombre completo: & \linead{10cm}\\
Fecha: & \linead{4cm} \quad Curso/grupo: \linead{4cm}\\
Mi firma: & \linead{9cm}\\
\end{tabularx}
\normalsize

\begin{infoband}{milcGradeDark}
\textbf{Para la próxima sesión.} Dos cosas. \textbf{Primero, ten la conversación} y trae \textbf{las cuatro respuestas} y \textbf{en qué te equivocaste} al imaginar el martes. \emph{Si puedes, pide algo más: acompañar a esa persona un rato en su trabajo. Media hora mirando enseña más que un semestre de folletos.} \textbf{Segundo, pregunta en tu casa por el oficio propio:} averigua con alguien mayor \textbf{cómo llegó al trabajo que hace o hizo} ---si lo escogió, si le tocó, si quería otra cosa---, \textbf{qué es lo que mejor sabe hacer} y \textbf{cómo sabe cuándo le quedó bien}. \textbf{Trae:} las respuestas de tu conversación y las de tu casa. \emph{Vas a encontrar dos cosas que suelen convivir: casi nadie de las generaciones anteriores escogió su oficio como se supone que ustedes deben escoger ---les tocó, o se fue dando---, y aun así casi todos tienen un criterio finísimo sobre cuándo su trabajo quedó bien hecho. Eso segundo no vino de haber escogido bien: vino de haberlo hecho muchas veces.}
\end{infoband}

\end{document}
